\documentclass[12pt]{article}
\usepackage{xeCJK}

\usepackage{charter}
\usepackage{fullpage}
\usepackage[colorlinks=false]{hyperref}
\usepackage{ifthen}
\usepackage{comment}
\usepackage[title,titletoc]{appendix}
\usepackage{pagecolor}
\usepackage{amsmath}
\usepackage{amsfonts}
%\usepackage[normalem]{ulem}
\usepackage{siunitx}
\usepackage{amsthm}
\sisetup{per=slash, load=abbr}
\usepackage{paralist}
\usepackage{pgfplots}
\usetikzlibrary{positioning}
\usetikzlibrary{fit}
\usetikzlibrary{snakes}
\usetikzlibrary{shapes.geometric}
\usetikzlibrary{patterns}
\usetikzlibrary{shapes,arrows,chains}
\usepgfplotslibrary{patchplots,colormaps}
\usetikzlibrary{calc}
\usetikzlibrary{positioning, fit}
\usetikzlibrary{backgrounds}
\usetikzlibrary{intersections}

\newcommand{\whitepaper}[1]{\begin{center}\fbox{\parbox{0.75\textwidth}{{\small
#1}}}\end{center}}

\newcommand{\pcolor}{white!25}

\usepackage{setspace}
\usepackage{algorithm2e}
\bibliographystyle{ieeetr}

\usepackage{geometry}
\geometry{left=3cm,right=3cm,top=1.6cm,bottom=3cm,headheight=0pt,headsep=1.5em}
\usepackage{fancyhdr}
\pagestyle{fancy}
%\chead{\includegraphics[scale=0.1]{../common/logo.png}}  %在此处插入logo.pdf图片 图片靠左
\chead{forum.gotoken.io}  %在此处插入logo.pdf图片 图片靠左
\lhead{} % 页眉中间位置内容
\rhead{}
\usepackage{expl3}
\ExplSyntaxOn
\newcommand\latinabbrev[1]{
  \peek_meaning:NTF . {% Same as \@ifnextchar
    #1\@}%
  { \peek_catcode:NTF a {% Check whether next char has same catcode as \'a, i.e., is a letter
      #1.\@ }%
    {#1.\@}}}
\ExplSyntaxOff

%Omit final dot from each def.
\def\eg{\latinabbrev{e.g}}
\def\etal{\latinabbrev{et al}}
\def\etc{\latinabbrev{etc}}
\def\ie{\latinabbrev{i.e}}
\newcommand{\dapp}{DApp\xspace}

%\setlength{\topskip}{1em}

\usepackage{indentfirst}
\usepackage{verbatimbox}
\usepackage{listings}

\lstset{basicstyle=\ttfamily}


\setCJKmainfont[BoldFont = STSongti-SC-Bold]{STSongti-SC-Regular}
\setCJKfamilyfont{hei}{SIL-Hei-Med-Jian}		%宋体
\setCJKfamilyfont{song}{SimSun}		%宋体
\setCJKfamilyfont{kai}{Kaiti}		%楷体
\setCJKfamilyfont{fang}{song}	%仿宋
\setCJKfamilyfont{li}{song}			%隶书
\setCJKfamilyfont{you}{Yuanti}		%幼圆

\newcommand{\song}{\CJKfamily{song}}	%宋体
\newcommand{\hei}{\CJKfamily{hei}}	%黑体
\newcommand{\kai}{\CJKfamily{kai}}	%楷体
\newcommand{\fs}{\CJKfamily{fang}}	%仿宋
\newcommand{\li}{\CJKfamily{li}}		%隶书
\newcommand{\you}{\CJKfamily{you}}	%幼圆
\newcommand{\reffig}[1]{图\ref{#1}}
\newcommand{\refsec}[1]{\S~\ref{#1}}

\newcommand{\txhash}[1]{{\href{https://etherscan.io/tx/#1}{\texttt{#1}}}}
\newcommand{\ethaddr}[1]{{\href{https://etherscan.io/address/#1}{\texttt{#1}}}}

\onehalfspacing   % ----------设置1.5倍行距（可能有意义，待调整）

%\parindent=20pt  % -------------------首行缩进大小，英文分段就直接0pt了吧。
\setlength{\parindent}{2.1em}
\setlength{\parskip}{0.3\baselineskip}
\newcommand{\nrcore}{Core Nebulas Rank}
\newcommand{\nrext}{Extended Nebulas Ranks}
\newcommand{\dom}{{\; \texttt{dom}\;}}
\newcommand{\nr}{\Gamma}
\newcommand{\ccontract}[3]{%
\begin{minipage}{0.9\linewidth}%
\centering
\begin{tikzpicture}[framed]%
\node[minimum width=0.9\textwidth, text width=0.9\textwidth, align=left] (a) {\ethaddr{#1}};%
\node[anchor=north west] (b) at ($(a.south west) + (1, -0.1)$) {\texttt{.#2(}};%
\node[anchor=north west] at ($(b.south west) + (1, -0.1)$) {\begin{minipage}{0.7\textwidth}%
{\texttt{#3)}}%
\end{minipage}};%
\end{tikzpicture}%
\end{minipage}%
}

\setCJKsansfont[BoldFont = STHeitiSC-Medium]{STHeitiSC-Light}


\newtheorem{property}{特征}
\newtheorem{corollary}{推论}
%\addbibresource{reference.bib}

\begin{document}
\pagestyle{empty}
\renewcommand{\contentsname}{目录}
\renewcommand{\abstractname}{摘要}
\renewcommand{\refname}{参考文献}
%\renewcommand{\nomname}{术语表（按首字母排序）}
\renewcommand{\figurename}{图}
\renewcommand{\tablename}{表}
\renewcommand{\baselinestretch}{1.5}
\renewcommand{\appendixname}{附录}
\renewcommand{\proofname}{证明}

\pagecolor{\pcolor}

\begin{titlepage}
  \begin{center}
    \vspace*{5.5cm}
    %\includegraphics[scale=0.5]{../common/logo.png}

    \vspace{0.5cm}


   \textbf{\huge{GoToken合约及操作记录}}

    \vspace{0.5cm}
    
    \textbf{}
  \end{center}

\end{titlepage}
\setcounter{page}{0}
%\thispagestyle{empty}
\tableofcontents
\newpage
\setcounter{page}{1}
\pagestyle{fancy}
\vspace*{0.01cm}

\section{介绍}
本文记录了GoToken合约情况，以及多签操作记录。个人的操作行为不在本文记录的范围内。

\section{工厂合约及地址}
\begin{enumerate}

\item \ethaddr{0x38c931c4064459468d4F6ED0Ae16F4Ba433Aa03F}，多签工厂合约，用于生成多签合约，{\color{red} 已废弃}。
\item \ethaddr{0xe6D887f548b48bF30b06a68a0aaE06F62A525032}，多签工厂合约，用于生成多签合约。
\item \ethaddr{0x98D5EC13d08E5796De7191c1719D6D5bDec4a74c}，可信列表工厂合约，用于生成一个可信的列表。
\item \ethaddr{0x5004e5C6b2D2009eB273bBaB45E5a756e2f729E5}，ERC20代币工厂合约，用于生成一个ERC20代币。
\item \ethaddr{0xfcD4bfFd7B76092c59Fb9791ff3F9B5405D4010A}，投资及分发工厂合约，业务相关。
\item \ethaddr{0x5b0D747aF4AF70d8022b1Cbe4Fa27010B3C3AD6d}，初始锁仓工厂合约，用于生成一个锁仓ERC20的合约。
\item \ethaddr{0x649C0bA9D1454fA22331B2e801505a234504B8BE}，GT项目工厂合约，用于生成一个GT项目的合约。
\end{enumerate}

\section{实例合约及地址}
{\color{red} 注意，本节仅描述了构造时使用的参数，因此随着后面的操作，本节描述的地址可能发生变化，以链上地址为准}

GoToken的合约由三个合约组成，分别为
\begin{enumerate}
\item \ethaddr{0x9FDB24Df185b4e6c42846c2f1355cA0A2BB7e043}，多签合约，用于管理GTToken合约的信任合约列表；{\color{red} 已变更为\ethaddr{0x8e61F120beAe3CF91cC9a017AC48844348e3b34A}}，构造参数为
\begin{verbatim}
(['0xe855B4cb17eA22CAE1be5FeB7fCDC0Ef67DCa84D',
  '0x3e6F107Fd4A95AF86108c1F44E502A6136AD386e',
  '0x57955d7AA271DbDDE92D67e0EF52D90c6E4089cA'])
\end{verbatim}
\item \ethaddr{0xBEb0FAE8c75c79e5f92c1C23C6435bD0509c276e}，GTToken的信任合约列表，使用多签操作；构造参数为
\begin{verbatim}
([], 
 0x9FDB24Df185b4e6c42846c2f1355cA0A2BB7e043) //多签地址, 已变更
\end{verbatim}
\item \ethaddr{0x353214343Ee192AD8a58C62961B972F4d5a6877E}，GTToken合约，这是一个标准的ERC20合约，用来完成代币的管理、转账等操作；构造参数为
\begin{verbatim}
('0x0000000000000000000000000000000000000000', 
 0, "GoToken", 6, "GOO", true, 
 0x9FDB24Df185b4e6c42846c2f1355cA0A2BB7e043, //多签地址, 已变更
 0xBEb0FAE8c75c79e5f92c1C23C6435bD0509c276e) //信任列表地址
\end{verbatim}
\end{enumerate}

GTToken合约的增发及销毁必须通过信任合约列表中的合约进行，目前在信任列表中存在着两个合约，
\begin{enumerate}
\item \ethaddr{0xAA9a51f48834924B2F79639Af74FefE9BFF74529} FundAndDistribute合约，用于完成铸币及代币分发操作；构造参数为
\begin{verbatim}
(0x353214343Ee192AD8a58C62961B972F4d5a6877E, //GTToken地址
 "USDT for GOO", "Only for Funders", 
 0xdAC17F958D2ee523a2206206994597C13D831ec7,  //USDT地址  
 0x9FDB24Df185b4e6c42846c2f1355cA0A2BB7e043,  //多签地址, 已变更
 0x735451974A28f63b0593cCe91696659c1B900380)  //信任列表地址
\end{verbatim}
注意，此处的信任地址列表为投资人白名单，虽然内容上与GTToken的信任合约列表一样，但是构造参数并不同，其构造参数为：
\begin{verbatim}
(['0xe855B4cb17eA22CAE1be5FeB7fCDC0Ef67DCa84D',
  '0x3e6F107Fd4A95AF86108c1F44E502A6136AD386e',
  '0x57955d7AA271DbDDE92D67e0EF52D90c6E4089cA'], //投资人列表
 0x9FDB24Df185b4e6c42846c2f1355cA0A2BB7e043) //多签地址
\end{verbatim}

\item \ethaddr{0x491c8a58DDb5a6380FeA7fE8909788f681453B00} InitLock合约，用于完成初始锁仓功能；构造参数为
\begin{verbatim}
(0x353214343Ee192AD8a58C62961B972F4d5a6877E, //GTToken地址
 9932664,  //解锁区块高度，大约锁仓6个月
 ['0xe855B4cb17eA22CAE1be5FeB7fCDC0Ef67DCa84D',
  '0x3e6F107Fd4A95AF86108c1F44E502A6136AD386e',
  '0x57955d7AA271DbDDE92D67e0EF52D90c6E4089cA'], //解锁分配地址
 [30000000000, 10000000000, 10000000000],   //解锁分配数量
 0x9FDB24Df185b4e6c42846c2f1355cA0A2BB7e043) //多签地址
\end{verbatim}
\end{enumerate}

不难发现，上述所有合约使用的多签合约为同一个合约地址，这仅仅是为了方便初期的管理，并不是为了限制。每个合约都可以实用\texttt{transfer\_multisig}方法修改其对应的多签合约。
\section{操作记录}
\paragraph{2019-10-28}
\begin{itemize}
\item \begin{verbatim}
add_multi_trusted(
  1,
  ['0xAA9a51f48834924B2F79639Af74FefE9BFF74529', //FundAndDistribute合约
   '0x491c8a58DDb5a6380FeA7fE8909788f681453B00'  //InitLock合约
  ])
\end{verbatim}
在GTToken合约中增加两个信任合约。
\item \begin{verbatim}
transfer(
  1, 
  0x3e6F107Fd4A95AF86108c1F44E502A6136AD386e,
  1000000)
\end{verbatim}
实验性的将一个GoToken打出（已经在事后打回并销毁）。
\end{itemize}

\paragraph{2019-11-17}
\begin{itemize}
\item 对新的多签工厂合约调用，\\\ethaddr{0xe6D887f548b48bF30b06a68a0aaE06F62A525032}
\begin{verbatim}
createMultiSig(['0xe855B4cb17eA22CAE1be5FeB7fCDC0Ef67DCa84D',
  '0x3e6F107Fd4A95AF86108c1F44E502A6136AD386e',
  '0x57955d7AA271DbDDE92D67e0EF52D90c6E4089cA'])
\end{verbatim}
返回地址为\ethaddr{0x8e61F120beAe3CF91cC9a017AC48844348e3b34A}，为GT系统新的多签合约；
\\\txhash{0x31288e10fc62a675c76e864b09b72146f156387ab60a990215c1260b81db498e}

\item 修改GTToken的信任列表合约的多签地址，\\\ethaddr{0xBEb0FAE8c75c79e5f92c1C23C6435bD0509c276e}
\begin{verbatim}
transfer_multisig(1, 0x8e61F120beAe3CF91cC9a017AC48844348e3b34A)
\end{verbatim}

\item 修改GTToken 合约的多签地址，\\\ethaddr{0x353214343Ee192AD8a58C62961B972F4d5a6877E}
\begin{verbatim}
transfer_multisig(2, 0x8e61F120beAe3CF91cC9a017AC48844348e3b34A)
\end{verbatim}

\item 修改FundAndDistribute 合约的多签地址，\\\ethaddr{0xAA9a51f48834924B2F79639Af74FefE9BFF74529}
\begin{verbatim}
transfer_multisig(3, 0x8e61F120beAe3CF91cC9a017AC48844348e3b34A)
\end{verbatim}

\item 修改FundAndDistribute 白名单合约的多签地址，\\\ethaddr{0x735451974A28f63b0593cCe91696659c1B900380}
\begin{verbatim}
transfer_multisig(4, 0x8e61F120beAe3CF91cC9a017AC48844348e3b34A)
\end{verbatim}

\end{itemize}

\paragraph{2019-11-28}
\begin{itemize}
\item 创建新的多签合约，用于管理FundAndDistribute白名单，与其他多签合约分离，\\
\ccontract{0xe6D887f548b48bF30b06a68a0aaE06F62A525032}{createMultiSig}{['0xe855B4cb17eA22CAE1be5FeB7fCDC0Ef67DCa84D',
  \\'0x3e6F107Fd4A95AF86108c1F44E502A6136AD386e',
  \\'0x57955d7AA271DbDDE92D67e0EF52D90c6E4089cA',
  \\'0x0aCA6913289d7887983F26d4CBF7ACF4B104D03f']}\\
返回的合约地址为\ethaddr{0x86880095C6A4Dccef7e8dBdB8A864d417A37b7EA};
\item 修改FundAndDistribute 白名单合约的多签地址，\\
\ccontract{0x735451974A28f63b0593cCe91696659c1B900380}{transfer\_multisig}{1, \\0x86880095C6A4Dccef7e8dBdB8A864d417A37b7EA}

\item 在FundAndDistribute白名单中增加地址，\\
\ccontract{0x735451974A28f63b0593cCe91696659c1B900380}{add\_trusted}{1, 
     \\0x0aCA6913289d7887983F26d4CBF7ACF4B104D03f}

\item 修改FundAndDistribute的多签地址，\\
\ccontract{0xAA9a51f48834924B2F79639Af74FefE9BFF74529}{transfer\_multisig}{2, \\0x86880095C6A4Dccef7e8dBdB8A864d417A37b7EA}

\end{itemize}


\newpage
\begin{appendices}
\section{操作帮助}
本节对于常见的操作给出通用的合约操作方法，本节描述的方法仅针对直接操作合约，以EtherScan为例。
\subsection{查看多签时使用的\texttt{id}}
在进行多签操作时，通常需要一个\texttt{id}，此\texttt{id}用于标识试图通过达成的共识的编号，该编号有以下特点：
\begin{enumerate}
\item 该\texttt{id}必须是从1开始连续递增的，不能跳过编号；
\item 一个多签人仅能对同一个\texttt{id}进行一次表决；
\item 该\texttt{id}的编号对于不同的函数是独立的，互相不影响；
\item 该\texttt{id}的编号是由多签合约提供的，因此多签合约更换时，相应的\texttt{id}需要重新开始编号。
\end{enumerate}
综合上述特征，不难发现，在多签操作前，取得正确的\texttt{id}是很重要的。

对于任一非多签合约，应该都可以在EtherScan中找到其对应的多签合约（合约中的\texttt{multisig\_contract}地址），
在多签合约中，调用(read contract)\texttt{get\_unused\_invoke\_id}，填入要进行多签的函数名，就可以得到相应的\texttt{id}。

\subsection{如何进行铸币}
铸币的前提是铸币人需要在铸币白名单中。

铸币操作分为两步：
\begin{enumerate}
\item 调用USDT合约的\texttt{approve}方法，需要注意的是，其中的spender必须为FundAndDistribute合约的地址，数量必须大于需要铸币的数量，例如\\
\ccontract{0xdAC17F958D2ee523a2206206994597C13D831ec7}{approve}{0xAA9a51f48834924B2F79639Af74FefE9BFF74529, \\10000000000} \\
其中，调用的地址为USDT合约的地址。

\item 调用FundAndDistribute合约的\texttt{fund}方法，需要注意的是，由于该方法需要进行多次转账，因此需要适当调高gas费用。
\end{enumerate}
%\input{changelog.tex}
\end{appendices}

\end{document}
